\section{Artisan of the Soul}

\begin{quotex}
All men and women are entrusted with the task of crafting their own life: in a certain sense, they are to make of it a work of art, a masterpiece.
\flright{\textsc{Pope John Paul II}, \emph{Letter to Artists}\footnote{\url{http://www.vatican.va/content/john-paul-ii/en/letters/1999/documents/hf_jp-ii_let_23041999_artists.html}}}
\end{quotex}


Your life is your soul. Like a bodybuilder caring for his physique, give your soul exercise and watch what it consumes. Food for the soul includes the images and thoughts you entertain. Be pure, be watchful, be wise.

\paragraph{Grievances}


\begin{quotex}
From this degenerate crop comes the revolutionary proletariat, with its hatred born of grievances, and the drawing-room Bolshevism of the aesthetes and literary folk, who enjoy and advertise the attractiveness of such states of mind.
\flright{\textsc{Oswald Spengler}, \emph{The Hour of Decision}}

The sense that one is being continually affronted is precisely a plebeian feeling.
\flright{\textsc{Nicolai Berdyaev}}
\end{quotex}


A world in which everyone takes offense, or claims to be a victim, is one ruled by the plebeian spirit and the pseudo-elite who encourage it.

\paragraph{The Quest for Order}


\begin{quotex}
The great man is filled with a very different passion, the will toward \emph{order}.
\flright{\textsc{Ezra Pound}}
\end{quotex}


The plebeian, on the other hand, is consumed with a passion for power.

\paragraph{Do It Like an Animal}


Several years ago I heard a faux bishop on TV claim that certain practices, which in the olden days were quaintly called ``unnatural acts'', could not be unnatural since animals in the wild performed them. Of course, in the olden days, bishops encouraged us to imitate the Saints, or even Christ, but never animals.

\paragraph{Affinity and Attraction}


The Law of Affinity, which governs births, has been debased by the New Agers into the ``Law of Attraction''. The former is an act of will and takes privation into account. The Law of Attraction, on the other hand, is more like magical thinking\footnote{\url{https://en.wikipedia.org/wiki/Magical_thinking}}: just by thinking of something, the impersonal universe will somehow provide. The persistence of this notion is fascinating, since prosperity, love, health, and so on, remain elusive.

Of course, the obvious corollary, therefore, is that the lack of such goods is attributable to ``stinking thinking.'' Several months ago, a distraught woman approached me for advice. Not only did her boyfriend run off with a younger woman, but he managed to take a lot of her cash and even her car. Rather than lending a sympathetic ear, her new age friends insisted that something in her mind had attracted such turmoil. In other words, she was to blame.

I asked her how she could be responsible for the free acts of another person. As for her other problems, I explained some aspects of civil law and suggested some lawyers. That gave her some peace of mind as well as a course of action. It wasn't much, but it was certainly better than having her meditate on her car coming back to her.

On the other hand, there is a certain basis to that perspective, despite its distortions. By the Law of Affinity, you have chosen, to some extent, those you come into contact with. That is worth pondering. So, if you did get involved in a bad relationship, look at it more closely. Were there little clues about that person's values or psychological problems? Was there some neediness in you that kept you in that relationship?

Ignorance and lack of consciousness are privations. The purpose of the exercise is not to place blame, but rather to lead to more personal power, or freedom.

\paragraph{New Sins}


Man is subjected to numerous laws at multiple levels. There are the abstract laws of logic and maths. There are physical laws of gravity and the other forces. It goes on: chemical, biological, sociological, and psychological laws all serve to limit one's freedom and power.

Then there are laws formulated by society. Positive law should be in conformance with natural law, or else it places unnecessary burdens. The modern world is not as liberating as it believes, since it has created a vast number of new laws, now defined as ``sins''. If it is not in Dante's hell or confession manuals published prior to 1789, it is not a sin.

\paragraph{Real Sins}


A real sin requires three characteristics:

\begin{itemize}


\item It must be a serious matter, i.e., a breach of cosmic order.

\item The person must know it to be so

\item The person must do it freely or willingly
\end{itemize}


The three conditions are not always met, even in acts we consider heinous. A person may not know or understand the cosmic order. More likely, he is hardly free, since he is placed under so many laws, not of his choosing.

Hence, the artisan of the soul will seek to know the cosmic order. He will work to become free from useless laws, psychological turmoil, or hastily made commitments. In that case, the experience of sin takes on a weightier, even terrifying, significance. The breach of the cosmic law involves a rupture with all creation.

Be careful about the opposites of the three imperatives.

\begin{itemize}


\item \textbf{Purity}: \emph{impurity, divided will, doubt}


\item \textbf{Watchfulness}: \emph{sleep, hypnotized}


\item \textbf{Wisdom}: \emph{slyness, cleverness}

\end{itemize}



\flrightit{Posted on 2015-03-05 by Cologero}

\begin{center}* * *\end{center}

\begin{footnotesize}
\begin{sffamily}
\texttt{Paulo on 2015-03-06 at 05:32 said:}

The question is

how to transcende

the laws???

\hfill

\texttt{Tom Blanchard on 2015-03-06 at 15:25 said:}

Interesting how the pursuit of order over power results in the acquisition of true power (freedom, as you note) as a natural consequence anyhow. As I believe Evola notes, power is feminine entity and naturally weds itself to those who are worthy of it.

Do you have access / reference to any kind of pre-1789 confession manual? I am actually making my first confession (decades of sins to recount) a week from now, so it's a bit top-of-mind for me.

(I will send Aeneas a bottle of Frangelico for each Evola text translated.)

\hfill

\texttt{Cologero on 2015-03-06 at 17:02 said:}

The section on the commandments should suffice, Mr Blanchard:

\textit{Catechism of Pope Pius X:} \url{http://www.gornahoor.net/library/CatechismSSPX.pdf}

\hfill

\texttt{Matthew Smallwood on 2015-03-07 at 10:26 said:}

Order is a dirty word these days, so it's interesting that you can blaspheme God in any way you want, but there are other words like ``order'' or ``sin'' which (if you say) get you branded as a heretic. ``Whoever is righteous has regard for the life of his beast, but the mercy of the wicked is cruel.'' Since men are living at the levels of beasts, things have to get much worse before we realize that the erstwhile rulers of the said modern world are not ``just'', although they may be ``socially just''. What continually makes me wonder is how the best \& brightest of Europe were seduced by classical liberalism, and almost to the man fell into the thought patterns that lead toward Revolution. In some cases, it is that the person communicating their insight has had their work distorted when passed on.

\hfill

\texttt{Paulo on 2015-03-07 at 11:34 said:}

Tom Blanchard

``power is feminine entity and naturally weds itself to those who are worthy of it.''

What a astonishing statement!?Shakti!This exaclty statement,it was yours or did you get in some place else?

See how good it is to find a hard woman?

\hfill

\texttt{Mark Citadel on 2015-04-03 at 16:25 said:}

Indeed Tom, one must be worthy of true power in order to receive it and own it. Our current elite operates on an illegitimate power that was never given to them. They stole it and abused it.

\end{sffamily}
\end{footnotesize}

